\documentclass[11pt]{article}

\usepackage[margin=1in]{geometry}
\usepackage[T1]{fontenc}
\usepackage[utf8]{inputenc}
\usepackage{lmodern}
\usepackage{amsmath, amssymb}
\usepackage{booktabs}
\usepackage{enumitem}
\usepackage{xcolor}
\usepackage{hyperref}
\usepackage{listings}
\usepackage{float}
\usepackage{titlesec}
\usepackage{tcolorbox}
\usepackage{array}
\usepackage{tabularx}
\usepackage{multirow}

\tcbuselibrary{skins, breakable}

%% ── colour palette ──────────────────────────────────────────────────────────
\definecolor{inquiryblue}{RGB}{20, 60, 120}
\definecolor{inquirybg}{RGB}{235, 242, 255}
\definecolor{warnorange}{RGB}{180, 80, 0}
\definecolor{warnbg}{RGB}{255, 243, 224}
\definecolor{goodgreen}{RGB}{0, 90, 40}
\definecolor{goodbg}{RGB}{220, 245, 230}
\definecolor{codegray}{RGB}{245, 245, 245}
\definecolor{critred}{RGB}{140, 20, 20}
\definecolor{critbg}{RGB}{255, 235, 235}

\hypersetup{
  colorlinks=true,
  linkcolor=inquiryblue,
  urlcolor=inquiryblue,
  citecolor=inquiryblue
}

%% ── tcolorbox environments ───────────────────────────────────────────────────
\newtcolorbox{inquiry}[1]{
  colback=inquirybg, colframe=inquiryblue,
  fonttitle=\bfseries\color{white},
  title={#1},
  breakable, arc=2pt, boxrule=1pt
}

\newtcolorbox{warning}[1]{
  colback=warnbg, colframe=warnorange,
  fonttitle=\bfseries\color{white},
  title={#1},
  breakable, arc=2pt, boxrule=1pt
}

\newtcolorbox{goodbox}[1]{
  colback=goodbg, colframe=goodgreen,
  fonttitle=\bfseries\color{white},
  title={#1},
  breakable, arc=2pt, boxrule=1pt
}

\newtcolorbox{critique}[1]{
  colback=critbg, colframe=critred,
  fonttitle=\bfseries\color{white},
  title={#1},
  breakable, arc=2pt, boxrule=1pt
}

\newtcolorbox{apispec}[1]{
  colback=codegray, colframe=black!40,
  fonttitle=\bfseries,
  title={API: \texttt{#1}},
  breakable, arc=2pt, boxrule=0.6pt
}

%% ── code style ───────────────────────────────────────────────────────────────
\lstdefinestyle{cstyle}{
  language=C,
  inputencoding=utf8,
  extendedchars=true,
  literate={─}{{-}}1,
  basicstyle=\ttfamily\small,
  keywordstyle=\color{inquiryblue}\bfseries,
  commentstyle=\color{goodgreen}\itshape,
  stringstyle=\color{warnorange},
  showstringspaces=false,
  frame=single,
  breaklines=true,
  columns=fullflexible,
  backgroundcolor=\color{codegray},
  rulecolor=\color{black!30},
  xleftmargin=6pt,
  xrightmargin=6pt,
  numbers=left,
  numberstyle=\tiny\color{black!40},
  numbersep=8pt
}

\lstdefinestyle{shellstyle}{
  inputencoding=utf8,
  extendedchars=true,
  literate={─}{{-}}1,
  basicstyle=\ttfamily\small,
  backgroundcolor=\color{black!90},
  basicstyle=\ttfamily\small\color{white},
  frame=none, breaklines=true,
  xleftmargin=6pt, xrightmargin=6pt
}

%% ── section formatting ───────────────────────────────────────────────────────
\titleformat{\section}{\large\bfseries\color{inquiryblue}}{}{0em}{\thesection\quad}
\titleformat{\subsection}{\normalsize\bfseries\color{inquiryblue!80!black}}{}{0em}{\thesubsection\quad}

%% ── title ────────────────────────────────────────────────────────────────────
\title{
  \textbf{\large Phase 1: TSP Instance Loading \& Distance Matrix}\\[4pt]
  \large An Inquiry-Based Systems Programming Assignment\\[8pt]
  \normalsize\textit{Architecture First. Code Second.}
}
\author{}
\date{}

\begin{document}
\maketitle
\tableofcontents
\newpage

%% ─────────────────────────────────────────────────────────────────────────────
\section{Critique of the Two Proposal Documents}
%% ─────────────────────────────────────────────────────────────────────────────

Before introducing the revised assignment, it is worth auditing the two source documents against your stated learning goals. Your goals, ranked, are:

\begin{enumerate}[label=\textbf{\arabic*.}]
  \item \textbf{Primary:} System design --- understand architecture, analyse alternatives, weigh tradeoffs.
  \item \textbf{Secondary:} Learn to implement in C at a level that builds real familiarity.
  \item \textbf{Tertiary:} Arrive at a working, memory-safe, testable Phase 1 baseline.
\end{enumerate}

\subsection{Proposal A (Doc 1: ``Inquiry-Based Coding Assignment'')}

\begin{goodbox}{What Proposal A does well}
\begin{itemize}
  \item Strong on the \emph{why} behind every design decision. It explains \emph{why} flat arrays beat \texttt{double**}, not just which to use.
  \item Correctly identifies that constant-memory is \emph{not} a reliable target and gives exact arithmetic to prove it ($128 \times 128 \times 8 = 128$\,KB vs.\ 64\,KB budget).
  \item Provides a clean, explicit API with named functions, separating concerns properly.
  \item Offers extension paths for ambitious work without muddying the main path.
  \item Includes a concrete testing matrix (P-01 through M-01), which is essential for a verification mindset.
\end{itemize}
\end{goodbox}

\begin{critique}{What Proposal A fails to deliver for your goals}
\begin{itemize}
  \item \textbf{Functions are named but not specified.} You see \texttt{tsp\_instance\_load(const char *path, TSPInstance *out)} but nowhere is there a description of: what arguments mean, what return values mean, what the function guarantees on success, what it guarantees on failure, or what state \texttt{out} is left in when something goes wrong. For someone learning C, this is the exact information that needs to be explicit.
  \item \textbf{No hardware grounding.} It talks about flat arrays being good for GPU mirroring but does not explain \emph{why} at the memory-bus level. It mentions cache locality without explaining spatial vs.\ temporal locality, prefetching, or cache lines.
  \item \textbf{Inquiry boxes are vague.} ``Compare alternatives'' appears often, but the actual comparison is not scaffolded with enough concrete data to force a conclusion. A learner can read the inquiries and still not know how to decide.
  \item \textbf{No file-structure scaffolding.} A beginning systems programmer looking at this needs to know: where does each function live, what headers does each file include, and in what order should things be written. That guidance is absent.
\end{itemize}
\end{critique}

\subsection{Proposal B (Doc 2: ``Memory \& Parsing Audit'')}

\begin{goodbox}{What Proposal B does well}
\begin{itemize}
  \item Hardware framing is strong: it mentions syscall count, cache prefetchers, and \texttt{cudaMemcpy} count explicitly.
  \item The \texttt{sqrt} vs.\ \texttt{sqrtf} inquiry is excellent and points at a real C pitfall that most beginners miss.
  \item Modular naming (Module 1--4) makes the assignment feel manageable.
  \item The symmetry question at the end is a genuine design fork, not rhetorical.
\end{itemize}
\end{goodbox}

\begin{critique}{What Proposal B fails to deliver for your goals}
\begin{itemize}
  \item \textbf{The constant-memory analysis is exactly backwards.} Proposal B presents \texttt{float} as the solution to the 64\,KB constraint for 128 cities. Let's check: $128 \times 128 \times 4 = 65{,}536$ bytes $= 64$\,KB exactly. That is the \emph{nominal hardware limit}, not comfortable margin. Proposal A correctly identifies this. Designing a system where you pray that the hardware limit is not 63.5\,KB in practice is not architecture---it is wishful thinking.
  \item \textbf{The \texttt{uint16\_t} option is raised and never resolved.} Integer weights for Euclidean coordinates with decimal values require explicit rounding decisions (floor? round? TSPLIB rounding?) that have non-trivial effects on tour quality. Raising this option without providing the intellectual tools to evaluate it is harmful.
  \item \textbf{Function signatures are not given at all.} The assignment asks you to implement \texttt{tsp\_instance\_load}, \texttt{tsp\_instance\_free}, and \texttt{tsp\_build\_distance\_matrix} but does not specify return types, argument meanings, error semantics, or invariants.
  \item \textbf{The modular structure is not connected.} Modules 1--4 read like four separate assignments. There is no scaffolding that shows how the pieces connect: what struct one module produces, which pointer the next module consumes.
  \item \textbf{No test structure is given.} ``Write a test runner'' is not enough guidance for a C learner to understand how to write a minimal unit test in plain C without a framework.
\end{itemize}
\end{critique}

\begin{inquiry}{Synthesis verdict}
Neither proposal alone is sufficient for your goals. Proposal A has better design reasoning; Proposal B has better hardware grounding. The revised assignment below takes the strongest elements of both, corrects the constant-memory mistake, and adds the one thing both are missing: \textbf{complete function contracts that specify inputs, outputs, and behavioral guarantees in enough detail to write the code from.}
\end{inquiry}

\newpage

%% ─────────────────────────────────────────────────────────────────────────────
\section{The Revised Assignment: Phase 1 Foundations}
%% ─────────────────────────────────────────────────────────────────────────────

\subsection{What you are building}

You are implementing a data-loading and matrix-construction layer for a Traveling Salesman Problem (TSP) solver in C99. This layer sits at the foundation of a system that will later run a Genetic Algorithm on the CPU, then migrate the hottest kernel to CUDA.

The system pipeline for Phase 1 is:

\begin{center}
\ttfamily
\small
[.tsp file on disk] $\longrightarrow$ \textbf{parser} $\longrightarrow$ \textbf{coordinate arrays} $\longrightarrow$ \textbf{distance matrix} $\longrightarrow$ [solver, later]
\end{center}

\medskip

Phase 1 delivers exactly three things:
\begin{enumerate}[label=\arabic*.]
  \item A struct that holds a TSP instance in memory.
  \item A function that parses a TSPLIB-format file into that struct.
  \item A function that fills the distance matrix from the parsed coordinates.
\end{enumerate}

\subsection{File layout you will create}

\begin{lstlisting}[style=cstyle, language={}]
tsp/
  include/
    instance.h        <- struct definition, macro, function declarations
  src/
    instance.c        <- parsing, matrix construction, cleanup
  tests/
    test_instance.c   <- unit tests (no framework needed)
    fixtures/
      square_4.tsp              <- 4-city known geometry test
      malformed_bad_coord.tsp   <- line with letters instead of numbers
      malformed_missing_section.tsp  <- no NODE_COORD_SECTION
  Makefile
\end{lstlisting}

You will implement in this order: struct $\to$ cleanup $\to$ parser $\to$ matrix $\to$ tests.

\newpage

%% ─────────────────────────────────────────────────────────────────────────────
\section{Design Decision 1: Matrix Storage Layout}
%% ─────────────────────────────────────────────────────────────────────────────

This is the most fundamental architectural decision in Phase 1.

\subsection{The three options}

You are choosing how to represent a 2D grid of $n \times n$ distances in C memory. Here are the three serious candidates.

\textbf{Option A:} Array of pointers (\texttt{double **dist})
\begin{lstlisting}[style=cstyle]
double **dist = malloc(n * sizeof(double *));   // 1 malloc
for (int i = 0; i < n; i++)
    dist[i] = malloc(n * sizeof(double));       // n more mallocs
// Access: dist[i][j]
\end{lstlisting}

\textbf{Option B:} Flat 1D array (\texttt{double *dist})
\begin{lstlisting}[style=cstyle]
double *dist = malloc(n * n * sizeof(double));  // 1 malloc
// Access: dist[i * n + j]
// Or with a macro: DIST(inst, i, j)
\end{lstlisting}

\textbf{Option C:} No stored matrix --- compute on demand
\begin{lstlisting}[style=cstyle]
// No dist field at all.
// Call: sqrt((x[i]-x[j])^2 + (y[i]-y[j])^2) each time.
\end{lstlisting}

\subsection{Comparison table}

\begin{inquiry}{Work through this table before reading the recommendation}
Fill in each cell before reading further. Forced thinking before revealed answer.
\end{inquiry}

\begin{center}
\renewcommand{\arraystretch}{1.6}
\begin{tabularx}{\textwidth}{>{\bfseries}p{4.2cm} X X X}
\toprule
Property & \textbf{A: double**} & \textbf{B: flat double*} & \textbf{C: on demand} \\
\midrule
Number of \texttt{malloc} calls & $n + 1$ & 1 & 0 \\
Memory is contiguous? & No & Yes & N/A \\
Rows adjacent in memory? & No & Yes & N/A \\
\texttt{free()} complexity & loop + outer free & single free & nothing \\
Copy to GPU (\texttt{cudaMemcpy}) & $n$ calls & 1 call & N/A \\
Access syntax & \texttt{dist[i][j]} & \texttt{dist[i*n+j]} & function call \\
Spatial cache locality & Poor & Excellent & N/A \\
Work per distance query & O(1) & O(1) & 1 sqrt + arithmetic \\
Repeated access cost & Equal & Equal & Each call is paid again \\
\bottomrule
\end{tabularx}
\end{center}

\subsection{The hardware reason flat arrays win}

When the CPU reads memory, it does not load one byte---it loads a \textbf{cache line} (typically 64 bytes on modern hardware). If your data is contiguous, loading one element brings its neighbours into cache automatically. This is called \emph{spatial locality}, and it means the hardware prefetcher can predict your next access and load it before you ask.

With \texttt{double**}, each row is a separate allocation somewhere in the heap. Row $i$ and row $i+1$ are not adjacent. When you finish row $i$ and start row $i+1$, the prefetcher cannot help you. This causes a cache miss for every row boundary.

For GPU transfer: \texttt{cudaMemcpy} expects a single contiguous source address and a byte count. A flat array gives you that in one call. A \texttt{double**} requires either $n$ individual copies (expensive overhead) or a reformatting step before transfer.

\begin{goodbox}{Decision: Use Option B}
Flat 1D array, accessed with a macro. Reason: single allocation, contiguous memory, one-call GPU copy, simple ownership model. Option A is the beginner trap in C. Option C is only interesting if memory is extremely constrained and access is sparse---not true for a dense GA solver.
\end{goodbox}

\subsection{The access macro and why it works}

In C, a 2D index $(i, j)$ in a row-major flat array of width $n$ is computed as $i \cdot n + j$. You will expose this through a macro so that code that uses the matrix does not need to know the internal layout:

\begin{lstlisting}[style=cstyle]
#define DIST(inst, i, j)  ((inst)->dist[(i) * (inst)->n + (j)])
\end{lstlisting}

The extra parentheses around every argument are mandatory. Without them, a call like \texttt{DIST(inst, a+1, b)} would expand to \texttt{inst->dist[a+1 * n + b]}, which is wrong due to operator precedence. Parenthesise macro arguments defensively, always.

\newpage

%% ─────────────────────────────────────────────────────────────────────────────
\section{Design Decision 2: Distance Matrix Precision}
%% ─────────────────────────────────────────────────────────────────────────────

\subsection{The type candidates}

\begin{center}
\renewcommand{\arraystretch}{1.5}
\begin{tabular}{lccccl}
\toprule
Type & Bytes & Range & Decimal digits & $128^2$ footprint & Notes \\
\midrule
\texttt{double} & 8 & $\pm 1.8 \times 10^{308}$ & $\approx 15$--16 & 128 KB & IEEE 754 double precision \\
\texttt{float}  & 4 & $\pm 3.4 \times 10^{38}$  & $\approx 6$--7  &  64 KB & IEEE 754 single precision \\
\texttt{uint16\_t} & 2 & 0--65535 & integer only & 32 KB & Requires explicit rounding \\
\bottomrule
\end{tabular}
\end{center}

\subsection{The GPU constant-memory calculation (done correctly)}

Both source proposals discuss GPU \texttt{\_\_constant\_\_} memory. You should understand this, but do not let it drive your Phase 1 decision.

\textbf{Nominal constant memory budget on NVIDIA hardware:} 64 KB.

\begin{align*}
128 \times 128 \times 8 &= 131{,}072 \text{ bytes} = 128 \text{ KB} \quad \text{(double --- too large)} \\
128 \times 128 \times 4 &= 65{,}536 \text{ bytes} = 64 \text{ KB} \quad \text{(float --- exactly the limit)} \\
128 \times 128 \times 2 &= 32{,}768 \text{ bytes} = 32 \text{ KB} \quad \text{(uint16\_t --- fits comfortably)}
\end{align*}

\begin{warning}{The constant-memory trap}
Proposal B suggests using \texttt{float} because $128^2 \times 4 = 64$\,KB fits the constant memory budget. This reasoning has two problems:

\begin{enumerate}[label=\arabic*.]
  \item 64\,KB is the \emph{nominal} maximum. In practice, the CUDA runtime and driver also use constant memory for internal variables. Your available budget is typically less than 64\,KB.
  \item Constant memory is not a scratchpad you can always use. The GPU broadcasts from constant memory efficiently \emph{only when all threads in a warp access the same address simultaneously}. A GA kernel where each thread follows a different tour accesses distances non-uniformly, which means constant memory may actually be \emph{slower} than L2 cache for this application.
\end{enumerate}

The right design choice is: \textbf{use what is correct and clean in Phase 1}. We will decide GPU memory placement in Phase 3, when we have profiling data.
\end{warning}

\subsection{The \texttt{uint16\_t} option evaluated}

Proposal B raises \texttt{uint16\_t}. This is a valid option for solvers that use TSPLIB instances with \emph{integer weights}, but not for Euclidean instances. The reason:

Euclidean distance between $(0,0)$ and $(1,1)$ is $\sqrt{2} \approx 1.41421356\ldots$. To store this as \texttt{uint16\_t = 1} you must round. TSPLIB defines a specific rounding convention (\texttt{NINT}, which is nearest-integer). Using it changes the problem slightly. For Phase 1 you are building a general Euclidean solver, so discard \texttt{uint16\_t} unless you later decide to implement TSPLIB integer-weight mode.

\begin{goodbox}{Decision: Use \texttt{double} for Phase 1}
\begin{itemize}
  \item \texttt{double} is the natural type for Euclidean geometry in C. \texttt{sqrt} returns \texttt{double}; coordinates parsed with \texttt{\%lf} are \texttt{double}.
  \item No implicit casts, no rounding decisions, no precision surprises.
  \item The CPU baseline is for correctness first. Memory pressure optimisation is a Phase 3 concern.
  \item If you intentionally choose \texttt{float} for Phase 1, you must document that you understand the above and explain why the reduced precision is acceptable.
\end{itemize}
\end{goodbox}

\newpage

%% ─────────────────────────────────────────────────────────────────────────────
\section{Design Decision 3: Full Matrix vs.\ Packed Symmetric Storage}
%% ─────────────────────────────────────────────────────────────────────────────

The Euclidean distance matrix is symmetric: $D_{ij} = D_{ji}$. You can exploit this.

\subsection{Two representations}

\textbf{Full $n \times n$ matrix:}
\begin{itemize}
  \item Memory: $n^2$ entries.
  \item Index: \texttt{dist[i*n + j]} for any $(i, j)$.
  \item Diagonal: $D_{ii} = 0$ explicitly stored.
\end{itemize}

\textbf{Packed upper-triangular (only $j > i$):}
\begin{itemize}
  \item Memory: $n(n-1)/2$ entries.
  \item Index: requires a formula. For $i < j$: index $= i \cdot n - i(i+1)/2 + j - i - 1$.
  \item Diagonal: not stored; always implicitly 0.
\end{itemize}

\begin{center}
\renewcommand{\arraystretch}{1.5}
\begin{tabular}{lcc}
\toprule
Property & Full $n \times n$ & Packed upper triangle \\
\midrule
Memory for $n=128$, \texttt{double} & 128 KB & 63.5 KB \\
Memory for $n=512$, \texttt{double} & 2 MB & $\approx$ 1 MB \\
Index expression & \texttt{i*n+j} & $in - i(i+1)/2 + j - i - 1$ \\
Can use same DIST macro? & Yes & No; needs new macro \\
Access pattern for GA & Uniform & Uniform \\
Copy to GPU & Trivial & Trivial \\
Debuggability & Easy (print as grid) & Hard (packed blob) \\
\bottomrule
\end{tabular}
\end{center}

\begin{inquiry}{Before reading the decision, answer this}
If you pack only the upper triangle and the solver code calls \texttt{DIST(inst, 5, 2)} where $5 > 2$, what happens? How would you handle it?
\end{inquiry}

Answer: your index formula is only defined for $i < j$. A call with $i > j$ gives a wrong result silently. You must either (a) swap $i$ and $j$ in the macro when $i > j$, or (b) enforce at the call site that queries always use the canonical order. Both add complexity.

\begin{goodbox}{Decision: Use full $n \times n$ matrix for Phase 1}
Reason: simplicity, debuggability, and uniform access with no conditional in the index. The memory waste is real ($\approx 50\%$) but small in absolute terms at Phase 1 sizes. Pack only if profiling shows memory is a bottleneck, or if you implement an explicit extension.
\end{goodbox}

\newpage

%% ─────────────────────────────────────────────────────────────────────────────
\section{Design Decision 4: Parsing Strategy}
%% ─────────────────────────────────────────────────────────────────────────────

\subsection{The two parsing approaches}

\textbf{Option A: \texttt{fscanf} directly}
\begin{lstlisting}[style=cstyle]
int id;
double x, y;
while (fscanf(fp, "%d %lf %lf", &id, &x, &y) == 3) {
    // store x, y
}
\end{lstlisting}

\textbf{Option B: \texttt{fgets} + \texttt{sscanf}}
\begin{lstlisting}[style=cstyle]
char line[256];
while (fgets(line, sizeof(line), fp)) {
    int id;
    double x, y;
    if (sscanf(line, "%d %lf %lf", &id, &x, &y) == 3) {
        // store x, y
    } else {
        // handle malformed line
    }
}
\end{lstlisting}

\subsection{Comparison}

\begin{center}
\renewcommand{\arraystretch}{1.5}
\begin{tabular}{lcc}
\toprule
Property & \texttt{fscanf} & \texttt{fgets} + \texttt{sscanf} \\
\midrule
What if a line is 10,000 chars? & Reads until whitespace, no truncation risk & Truncates to buffer; detectable \\
What if a field is ``ABC''? & Returns 0 matches; position hard to recover & Line is already buffered; easy to skip \\
Error recovery position in file & Ambiguous & Clean: you are at line start \\
Can you re-examine a bad line? & Only with ungetc tricks & Yes: you have it in \texttt{line[]} \\
Malformed file: can you free before exit? & Harder (parse state unclear) & Easier (line-level granularity) \\
\bottomrule
\end{tabular}
\end{center}

\begin{goodbox}{Decision: Use \texttt{fgets} + \texttt{sscanf}}
The line-buffer approach keeps the parse state clear at all times. If a line fails to parse, you know exactly where you are in the file, you can log or reject cleanly, and you can call \texttt{tsp\_instance\_free} before returning an error. This is the correct choice for any parser that must handle malformed input gracefully in C.
\end{goodbox}

\subsection{One-pass vs.\ two-pass allocation}

When parsing coordinates, you need to allocate arrays. You do not know $n$ until you have scanned the file. Two strategies:

\textbf{Two-pass:} scan the file once counting cities, seek back to the start, allocate exactly, parse again.\\
\textbf{One-pass:} grow arrays dynamically with \texttt{realloc} as you read.

\begin{inquiry}{Two-pass or one-pass?}
\begin{enumerate}[label=\alph*.]
  \item What does \texttt{realloc} do if the new size is larger? What does it do to the old pointer?
  \item In a two-pass parser, what happens if the file is on a pipe or socket rather than disk?
  \item Which approach has more predictable memory access patterns during parsing?
  \item TSPLIB files often have a \texttt{DIMENSION} header field. If you parse it first, can you do single-pass allocation without \texttt{realloc}?
\end{enumerate}
\end{inquiry}

\textbf{Recommended for this assignment:} parse the \texttt{DIMENSION} header, allocate once with known size, then parse \texttt{NODE\_COORD\_SECTION} in one pass. This avoids both \texttt{realloc} complexity and a second file scan.

\newpage

%% ─────────────────────────────────────────────────────────────────────────────
\section{The Complete API Specification}
%% ─────────────────────────────────────────────────────────────────────────────

This section gives you complete function contracts. A \emph{contract} specifies: what you must provide (preconditions), what the function guarantees when it returns (postconditions), and what it guarantees on both success and failure paths.

\subsection{The data structure (\texttt{include/instance.h})}

\begin{lstlisting}[style=cstyle]
#ifndef TSP_INSTANCE_H
#define TSP_INSTANCE_H

/* ── TSPInstance ────────────────────────────────────────────────────────────
 *
 *  Holds a single TSP instance loaded from a TSPLIB coordinate file.
 *
 *  Fields:
 *    n         Number of cities. Always >= 1 when the struct is valid.
 *              Zero indicates an uninitialized or freed struct.
 *
 *    coords_x  Pointer to an array of n x-coordinates, allocated on the heap.
 *              coords_x[i] is the x-coordinate of city i (0-indexed).
 *              NULL if not yet allocated.
 *
 *    coords_y  Pointer to an array of n y-coordinates, allocated on the heap.
 *              coords_y[i] is the y-coordinate of city i (0-indexed).
 *              NULL if not yet allocated.
 *
 *    dist      Pointer to a flat row-major n*n distance matrix, heap-allocated.
 *              dist[i*n + j] is the Euclidean distance from city i to city j.
 *              NULL if not yet built.
 *
 *  Invariants (when valid):
 *    dist[i*n + i] == 0.0          for all i
 *    dist[i*n + j] == dist[j*n+i]  for all i, j  (symmetry)
 *    dist[i*n + j] >= 0.0          for all i, j
 *
 *  Ownership:
 *    The struct owns all three pointers. tsp_instance_free() must be called
 *    exactly once on any struct touched by tsp_instance_load(), whether or
 *    not the load succeeded.
 * ─────────────────────────────────────────────────────────────────────────── */
typedef struct {
    int     n;
    double *coords_x;
    double *coords_y;
    double *dist;
} TSPInstance;

/* Access macro. i and j must be in [0, inst->n). */
#define DIST(inst, i, j)  ((inst)->dist[(i) * (inst)->n + (j)])

/* ── Function declarations ───────────────────────────────────────────────── */

int  tsp_instance_load(const char *path, TSPInstance *out);
void tsp_instance_free(TSPInstance *inst);
int  tsp_build_distance_matrix(TSPInstance *inst);

#endif /* TSP_INSTANCE_H */
\end{lstlisting}

\newpage

\begin{apispec}{tsp\_instance\_load}

\textbf{Signature:}
\begin{lstlisting}[style=cstyle, numbers=none]
int tsp_instance_load(const char *path, TSPInstance *out);
\end{lstlisting}

\textbf{Arguments:}
\begin{itemize}
  \item \texttt{path} --- A null-terminated C string giving the file system path to a TSPLIB-format \texttt{.tsp} file. Must not be \texttt{NULL}.
  \item \texttt{out} --- Pointer to a \texttt{TSPInstance} that the caller provides. The function writes the loaded instance into this struct. Must not be \texttt{NULL}. The struct's contents before the call are ignored; the function overwrites them.
\end{itemize}

\textbf{Returns:}
\begin{itemize}
  \item \texttt{0} on success. \texttt{out} is fully populated: \texttt{n} $\geq$ 1, \texttt{coords\_x} and \texttt{coords\_y} are heap-allocated arrays of length \texttt{n}, \texttt{dist} is \texttt{NULL} (matrix is not built by this function).
  \item \texttt{-1} on any failure. Failure conditions include: file not found, file is empty, \texttt{NODE\_COORD\_SECTION} is missing, any coordinate line is malformed, any allocation fails. On failure, \texttt{out} is left in a state that is safe to pass to \texttt{tsp\_instance\_free()}: all pointers it owns are either valid or \texttt{NULL}.
\end{itemize}

\textbf{What it does (step by step):}
\begin{enumerate}[label=\arabic*.]
  \item Initialise \texttt{*out} to the zero struct (\texttt{n=0}, all pointers \texttt{NULL}). This ensures \texttt{free} is safe even if we return early.
  \item Open the file at \texttt{path}. Return \texttt{-1} if it cannot be opened.
  \item Scan lines until \texttt{DIMENSION: N} (or \texttt{DIMENSION : N}) is found. Parse $N$ as the city count. Return \texttt{-1} if not found before \texttt{NODE\_COORD\_SECTION} or EOF.
  \item Allocate \texttt{coords\_x[N]} and \texttt{coords\_y[N]}. Return \texttt{-1} if either allocation fails.
  \item Scan lines until \texttt{NODE\_COORD\_SECTION} is found. Return \texttt{-1} if not found.
  \item For each of the next $N$ lines, parse with \texttt{sscanf(line, "\%d \%lf \%lf", \&id, \&x, \&y)}. If any line does not produce exactly 3 fields, return \texttt{-1}.
  \item Store \texttt{x} in \texttt{coords\_x[id-1]} and \texttt{y} in \texttt{coords\_y[id-1]}. (TSPLIB IDs are 1-indexed; convert to 0-indexed.)
  \item Close the file. Set \texttt{out->n = N}. Return \texttt{0}.
\end{enumerate}

\textbf{What it does NOT do:}
\begin{itemize}
  \item It does \emph{not} build the distance matrix. Call \texttt{tsp\_build\_distance\_matrix} separately.
  \item It does \emph{not} validate that city IDs form a contiguous sequence starting at 1.
\end{itemize}

\end{apispec}

\vspace{8pt}

\begin{apispec}{tsp\_instance\_free}

\textbf{Signature:}
\begin{lstlisting}[style=cstyle, numbers=none]
void tsp_instance_free(TSPInstance *inst);
\end{lstlisting}

\textbf{Arguments:}
\begin{itemize}
  \item \texttt{inst} --- Pointer to a \texttt{TSPInstance}. May be \texttt{NULL} (the function is a no-op in that case). The struct it points to may be partially initialised (some pointers non-\texttt{NULL}, others \texttt{NULL}).
\end{itemize}

\textbf{Returns:} Nothing.

\textbf{What it does:}
\begin{enumerate}[label=\arabic*.]
  \item If \texttt{inst} is \texttt{NULL}, return immediately.
  \item Call \texttt{free(inst->coords\_x)} if it is not \texttt{NULL}. Set to \texttt{NULL}.
  \item Call \texttt{free(inst->coords\_y)} if it is not \texttt{NULL}. Set to \texttt{NULL}.
  \item Call \texttt{free(inst->dist)} if it is not \texttt{NULL}. Set to \texttt{NULL}.
  \item Set \texttt{inst->n = 0}.
\end{enumerate}

\textbf{Why nullify after free?} Calling \texttt{free()} on a pointer does not set the pointer to \texttt{NULL}. If anything later accidentally calls \texttt{free()} again on the same pointer, it is undefined behaviour (double-free). Setting to \texttt{NULL} after freeing makes the struct safe to call \texttt{free} on a second time (\texttt{free(NULL)} is always a no-op in C).

\textbf{Why safe on partial init?} During \texttt{tsp\_instance\_load}, several allocations happen. If the third allocation fails and you return \texttt{-1}, the struct has some non-\texttt{NULL} pointers and some \texttt{NULL} ones. If \texttt{free} only frees non-\texttt{NULL} pointers, the caller can always call \texttt{tsp\_instance\_free} unconditionally without leaking.

\end{apispec}

\vspace{8pt}

\begin{apispec}{tsp\_build\_distance\_matrix}

\textbf{Signature:}
\begin{lstlisting}[style=cstyle, numbers=none]
int tsp_build_distance_matrix(TSPInstance *inst);
\end{lstlisting}

\textbf{Arguments:}
\begin{itemize}
  \item \texttt{inst} --- Pointer to a \texttt{TSPInstance} that has already been successfully loaded: \texttt{inst->n} $\geq 1$, \texttt{coords\_x} and \texttt{coords\_y} are valid, \texttt{dist} is \texttt{NULL} (or the function will overwrite it).
\end{itemize}

\textbf{Returns:}
\begin{itemize}
  \item \texttt{0} on success. \texttt{inst->dist} points to a heap-allocated array of \texttt{n*n} doubles satisfying the invariants (diagonal zero, symmetric, non-negative).
  \item \texttt{-1} if \texttt{malloc} fails. \texttt{inst->dist} is left \texttt{NULL}.
\end{itemize}

\textbf{What it does:}
\begin{enumerate}[label=\arabic*.]
  \item Allocate \texttt{n*n} doubles. Return \texttt{-1} on failure.
  \item For all $i$: set \texttt{DIST(inst,i,i) = 0.0}.
  \item For all $i < j$: compute $d = \sqrt{(x_i - x_j)^2 + (y_i - y_j)^2}$ using \texttt{sqrt()} from \texttt{<math.h>}. Set both \texttt{DIST(inst,i,j) = d} and \texttt{DIST(inst,j,i) = d}.
  \item Return \texttt{0}.
\end{enumerate}

\textbf{Why only the upper triangle?} Computing only $j > i$ means each pair is computed once. Mirroring enforces $D_{ij} = D_{ji}$ \emph{by construction} from the same value, not from two independent computations that might differ by a floating-point epsilon.

\textbf{The \texttt{sqrt} vs \texttt{sqrtf} issue:} \texttt{sqrt()} from \texttt{<math.h>} takes and returns \texttt{double}. \texttt{sqrtf()} takes and returns \texttt{float}. Since both coordinates and distances are \texttt{double} in this design, use \texttt{sqrt()}. If you used \texttt{float} and called \texttt{sqrt()} instead of \texttt{sqrtf()}, the argument would be implicitly promoted to \texttt{double} for the computation and then truncated back to \texttt{float} on assignment---not undefined behaviour, but needless and slightly imprecise. Always match the function suffix to your type.

\end{apispec}

\newpage

%% ─────────────────────────────────────────────────────────────────────────────
\section{The TSPLIB File Format (Narrow Subset)}
%% ─────────────────────────────────────────────────────────────────────────────

You will support one strict subset of TSPLIB. Not all variants. This is intentional: a parser that handles one format correctly is better than a parser that handles ten formats inconsistently.

\subsection{Format you will support}

\begin{lstlisting}[style=cstyle, language={}, numbers=none]
NAME: square_4
TYPE: TSP
COMMENT: 4-city unit square
DIMENSION: 4
EDGE_WEIGHT_TYPE: EUC_2D
NODE_COORD_SECTION
1 0.0 0.0
2 1.0 0.0
3 1.0 1.0
4 0.0 1.0
EOF
\end{lstlisting}

\textbf{Rules your parser enforces:}
\begin{itemize}
  \item Header fields may appear in any order before \texttt{NODE\_COORD\_SECTION}.
  \item \texttt{DIMENSION} must appear and must be a positive integer.
  \item \texttt{NODE\_COORD\_SECTION} must appear.
  \item After \texttt{NODE\_COORD\_SECTION}, each line must parse as \texttt{int double double} (id, x, y).
  \item The parser reads exactly \texttt{DIMENSION} coordinate lines.
  \item \texttt{EOF} at end is optional; the parser stops when it has read all $n$ cities.
\end{itemize}

\textbf{What you do NOT need to handle (yet):}
\begin{itemize}
  \item \texttt{EDGE\_WEIGHT\_SECTION} (explicit weight matrices)
  \item \texttt{DISPLAY\_DATA\_SECTION}
  \item \texttt{EDGE\_WEIGHT\_TYPE} values other than \texttt{EUC\_2D}
  \item 3D coordinates
\end{itemize}

\subsection{Test fixture: \texttt{square\_4.tsp}}

Create this file exactly. You will use it to verify every geometric invariant.

\begin{lstlisting}[style=cstyle, language={}, numbers=none]
NAME: square_4
TYPE: TSP
DIMENSION: 4
EDGE_WEIGHT_TYPE: EUC_2D
NODE_COORD_SECTION
1 0.0 0.0
2 1.0 0.0
3 1.0 1.0
4 0.0 1.0
EOF
\end{lstlisting}

Cities 0--3 (0-indexed after loading) form a unit square. The expected distance matrix is:

\[
D = \begin{pmatrix}
0 & 1 & \sqrt{2} & 1 \\
1 & 0 & 1 & \sqrt{2} \\
\sqrt{2} & 1 & 0 & 1 \\
1 & \sqrt{2} & 1 & 0
\end{pmatrix}
\]

\newpage

%% ─────────────────────────────────────────────────────────────────────────────
\section{C Implementation Guide}
%% ─────────────────────────────────────────────────────────────────────────────

This section gives you enough C scaffolding that you can implement each function without guessing syntax, while still writing the logic yourself.

\subsection{Step 1: Struct, macro, and static initialiser}

In \texttt{include/instance.h}, place the struct and macro as shown in Section~6.

In \texttt{src/instance.c}, write a static helper to initialise a struct to the safe zero state:

\begin{lstlisting}[style=cstyle]
#include "instance.h"
#include <stdio.h>
#include <stdlib.h>
#include <string.h>
#include <math.h>

/* Zero-initialise: all pointers NULL, n = 0.
 * After this call, tsp_instance_free(inst) is always safe. */
static void tsp_instance_init(TSPInstance *inst)
{
    inst->n        = 0;
    inst->coords_x = NULL;
    inst->coords_y = NULL;
    inst->dist     = NULL;
}
\end{lstlisting}

Write this first. Every other function depends on this invariant.

\subsection{Step 2: Implement \texttt{tsp\_instance\_free}}

Before writing the parser, write cleanup. The reason is: every early-return error path in the parser will call \texttt{tsp\_instance\_free}. If cleanup is buggy, your error paths are buggy.

Pattern to follow:

\begin{lstlisting}[style=cstyle]
void tsp_instance_free(TSPInstance *inst)
{
    if (inst == NULL) return;
    free(inst->coords_x);  inst->coords_x = NULL;
    free(inst->coords_y);  inst->coords_y = NULL;
    free(inst->dist);      inst->dist     = NULL;
    inst->n = 0;
}
\end{lstlisting}

\begin{inquiry}{Why is \texttt{free(NULL)} safe?}
The C standard (C99 \S7.20.3.2) guarantees that \texttt{free(NULL)} has no effect. This means you never need to check for \texttt{NULL} before calling \texttt{free}. But you \emph{do} need to set pointers to \texttt{NULL} after freeing them, to prevent accidental double-free.
\end{inquiry}

\subsection{Step 3: Skeleton of \texttt{tsp\_instance\_load}}

\begin{lstlisting}[style=cstyle]
int tsp_instance_load(const char *path, TSPInstance *out)
{
    tsp_instance_init(out);   /* safe even if we return -1 early */

    FILE *fp = fopen(path, "r");
    if (fp == NULL) return -1;

    char line[512];
    int  dimension_found = 0;
    int  n = 0;

    /* ── Pass 1: scan headers for DIMENSION ─────────────────────── */
    /* YOUR CODE: loop with fgets, look for "DIMENSION" keyword,
       parse the integer, set dimension_found = 1 when found.
       Stop when you hit NODE_COORD_SECTION or EOF.               */

    if (!dimension_found) { fclose(fp); return -1; }

    /* ── Allocate coordinate arrays ──────────────────────────────── */
    out->coords_x = malloc(n * sizeof(double));
    out->coords_y = malloc(n * sizeof(double));
    if (!out->coords_x || !out->coords_y) {
        fclose(fp);
        tsp_instance_free(out);  /* free whatever was allocated */
        return -1;
    }

    /* ── Find NODE_COORD_SECTION ─────────────────────────────────── */
    /* YOUR CODE: continue scanning lines until NODE_COORD_SECTION.
       If not found before EOF, call tsp_instance_free and return -1. */

    /* ── Parse coordinate lines ──────────────────────────────────── */
    for (int k = 0; k < n; k++) {
        if (fgets(line, sizeof(line), fp) == NULL) {
            fclose(fp); tsp_instance_free(out); return -1;
        }
        int   id;
        double x, y;
        /* YOUR CODE: sscanf the line. Check return value == 3.
           Store: out->coords_x[id-1] = x; out->coords_y[id-1] = y;
           If sscanf fails, call tsp_instance_free and return -1.   */
    }

    fclose(fp);
    out->n = n;
    return 0;
}
\end{lstlisting}

\subsection{Step 4: Implement \texttt{tsp\_build\_distance\_matrix}}

\begin{lstlisting}[style=cstyle]
int tsp_build_distance_matrix(TSPInstance *inst)
{
    int n = inst->n;
    inst->dist = malloc((size_t)n * n * sizeof(double));
    if (inst->dist == NULL) return -1;

    for (int i = 0; i < n; i++) {
        DIST(inst, i, i) = 0.0;
        for (int j = i + 1; j < n; j++) {
            double dx = inst->coords_x[i] - inst->coords_x[j];
            double dy = inst->coords_y[i] - inst->coords_y[j];
            double d  = sqrt(dx*dx + dy*dy);
            DIST(inst, i, j) = d;
            DIST(inst, j, i) = d;
        }
    }
    return 0;
}
\end{lstlisting}

Note: \texttt{(size\_t)n * n} is important. If \texttt{n} is large enough that \texttt{n*n} overflows a 32-bit \texttt{int}, the multiplication is done as \texttt{int} and wraps before the cast. Casting one operand to \texttt{size\_t} first forces the multiplication to happen in 64-bit.

\subsection{Step 5: The Makefile}

\begin{lstlisting}[style=cstyle, language=make]
CC      = gcc
CFLAGS  = -Wall -Wextra -Werror -std=c11 -Iinclude
LDFLAGS = -lm

all: test_instance

test_instance: src/instance.c tests/test_instance.c
	$(CC) $(CFLAGS) $^ $(LDFLAGS) -o $@

clean:
	rm -f test_instance

.PHONY: all clean
\end{lstlisting}

\newpage

%% ─────────────────────────────────────────────────────────────────────────────
\section{Testing: Writing Plain-C Unit Tests}
%% ─────────────────────────────────────────────────────────────────────────────

You are not using a test framework. You are writing a test runner in plain C. Here is the minimal pattern:

\begin{lstlisting}[style=cstyle]
/* tests/test_instance.c */
#include <stdio.h>
#include <stdlib.h>
#include <math.h>
#include "instance.h"

/* ── Test helpers ──────────────────────────────────────────────────────────── */
static int tests_run    = 0;
static int tests_passed = 0;

#define CHECK(condition, msg)                                       \
    do {                                                            \
        tests_run++;                                                \
        if (condition) {                                            \
            tests_passed++;                                         \
        } else {                                                    \
            fprintf(stderr, "FAIL [%s:%d] %s\n",                   \
                    __FILE__, __LINE__, msg);                       \
        }                                                           \
    } while (0)

#define CHECK_NEAR(a, b, tol, msg)  CHECK(fabs((a)-(b)) < (tol), msg)

/* ── Tests ─────────────────────────────────────────────────────────────────── */

static void test_load_square(void)
{
    TSPInstance inst;
    int rc = tsp_instance_load("tests/fixtures/square_4.tsp", &inst);

    CHECK(rc == 0,        "P-01: load returns 0 on success");
    CHECK(inst.n == 4,    "P-02: city count is 4");

    if (rc != 0) { tsp_instance_free(&inst); return; }

    /* P-03: spot-check one coordinate */
    CHECK_NEAR(inst.coords_x[0], 0.0, 1e-9, "P-03: city 0 x == 0.0");
    CHECK_NEAR(inst.coords_y[0], 0.0, 1e-9, "P-03: city 0 y == 0.0");

    int mrc = tsp_build_distance_matrix(&inst);
    CHECK(mrc == 0, "matrix build returns 0");

    /* D-01: diagonal is zero */
    for (int i = 0; i < inst.n; i++)
        CHECK(DIST(&inst, i, i) == 0.0, "D-01: diagonal is zero");

    /* D-02: symmetry */
    for (int i = 0; i < inst.n; i++)
        for (int j = 0; j < inst.n; j++)
            CHECK_NEAR(DIST(&inst, i, j), DIST(&inst, j, i), 1e-12, "D-02: symmetric");

    /* D-03: adjacent sides == 1 */
    CHECK_NEAR(DIST(&inst, 0, 1), 1.0, 1e-9, "D-03: dist(0,1)==1");
    CHECK_NEAR(DIST(&inst, 1, 2), 1.0, 1e-9, "D-03: dist(1,2)==1");

    /* D-04: diagonal == sqrt(2) */
    CHECK_NEAR(DIST(&inst, 0, 2), sqrt(2.0), 1e-9, "D-04: dist(0,2)==sqrt(2)");
    CHECK_NEAR(DIST(&inst, 1, 3), sqrt(2.0), 1e-9, "D-04: dist(1,3)==sqrt(2)");

    tsp_instance_free(&inst);
    /* M-01: after free, pointers are NULL */
    CHECK(inst.coords_x == NULL, "M-01: coords_x NULL after free");
    CHECK(inst.dist      == NULL, "M-01: dist NULL after free");
}

static void test_malformed_bad_coord(void)
{
    TSPInstance inst;
    int rc = tsp_instance_load("tests/fixtures/malformed_bad_coord.tsp", &inst);
    CHECK(rc == -1, "F-01: malformed coord returns -1");
    /* Must be safe to free even after failure */
    tsp_instance_free(&inst);
    CHECK(inst.n == 0, "F-01: n is 0 after failed load + free");
}

static void test_missing_section(void)
{
    TSPInstance inst;
    int rc = tsp_instance_load("tests/fixtures/malformed_missing_section.tsp", &inst);
    CHECK(rc == -1, "F-02: missing section returns -1");
    tsp_instance_free(&inst);
}

/* ── Main ───────────────────────────────────────────────────────────────────── */
int main(void)
{
    test_load_square();
    test_malformed_bad_coord();
    test_missing_section();

    printf("\nResults: %d/%d tests passed\n", tests_passed, tests_run);
    return (tests_passed == tests_run) ? 0 : 1;
}
\end{lstlisting}

\subsection{Malformed fixture files}

\textbf{\texttt{malformed\_bad\_coord.tsp}:}
\begin{lstlisting}[style=cstyle, language={}, numbers=none]
NAME: bad_coord
DIMENSION: 2
NODE_COORD_SECTION
1 0.0 0.0
2 NOTANUMBER 1.0
EOF
\end{lstlisting}

\textbf{\texttt{malformed\_missing\_section.tsp}:}
\begin{lstlisting}[style=cstyle, language={}, numbers=none]
NAME: missing_section
DIMENSION: 2
EOF
\end{lstlisting}

\subsection{Running under Valgrind}

\begin{lstlisting}[style=shellstyle]
make
valgrind --leak-check=full --error-exitcode=1 ./test_instance
\end{lstlisting}

Valgrind should report:
\begin{lstlisting}[style=shellstyle]
==PID== HEAP SUMMARY:
==PID==     in use at exit: 0 bytes in 0 blocks
==PID==  LEAK SUMMARY:
==PID==     definitely lost: 0 bytes in 0 blocks
\end{lstlisting}

If you see leaks, trace them by adding \texttt{--track-origins=yes}.

\newpage

%% ─────────────────────────────────────────────────────────────────────────────
\section{Test Matrix: What You Are Verifying}
%% ─────────────────────────────────────────────────────────────────────────────

\begin{center}
\renewcommand{\arraystretch}{1.6}
\begin{tabularx}{\textwidth}{llX}
\toprule
\textbf{ID} & \textbf{Category} & \textbf{What it checks} \\
\midrule
P-01 & Parsing & \texttt{tsp\_instance\_load} returns 0 on a valid file \\
P-02 & Parsing & Loaded city count matches \texttt{DIMENSION} \\
P-03 & Parsing & At least one coordinate value matches expected \\
D-01 & Matrix  & \texttt{DIST(inst, i, i) == 0.0} for all $i$ \\
D-02 & Matrix  & \texttt{DIST(inst, i, j) == DIST(inst, j, i)} for all $i, j$ \\
D-03 & Matrix  & Adjacent square sides equal 1.0 within tolerance \\
D-04 & Matrix  & Square diagonals equal $\sqrt{2}$ within tolerance \\
F-01 & Failure & Malformed coordinate line returns -1 \\
F-02 & Failure & Missing \texttt{NODE\_COORD\_SECTION} returns -1 \\
M-01 & Memory  & All pointers are \texttt{NULL} after \texttt{tsp\_instance\_free} \\
M-02 & Memory  & Valgrind reports zero leaked bytes on success path \\
M-03 & Memory  & Valgrind reports zero leaked bytes on failure path \\
\bottomrule
\end{tabularx}
\end{center}

\newpage

%% ─────────────────────────────────────────────────────────────────────────────
\section{Questions You Must Answer in Writing}
%% ─────────────────────────────────────────────────────────────────────────────

Answer these in a \texttt{DESIGN\_NOTES.md} file or as block comments at the top of \texttt{src/instance.c}. These are not optional. The point is to force you to connect the code you wrote to the reasoning behind it.

\begin{enumerate}[label=\textbf{Q\arabic*.}, leftmargin=*, itemsep=6pt]
  \item Why does \texttt{tsp\_instance\_free} set pointers to \texttt{NULL} after calling \texttt{free}? What specific undefined behaviour does this prevent?

  \item Why is \texttt{tsp\_instance\_init} called at the very beginning of \texttt{tsp\_instance\_load}, before the file is opened?

  \item Why does \texttt{tsp\_build\_distance\_matrix} iterate only for $j > i$ and then assign both \texttt{DIST(inst,i,j)} and \texttt{DIST(inst,j,i)} from the same value, rather than computing all $n^2$ entries independently?

  \item If you had used \texttt{double **dist} instead of \texttt{double *dist}, how many \texttt{free()} calls would \texttt{tsp\_instance\_free} need to make? Write out the loop.

  \item A colleague proposes: ``We should use \texttt{float} for the distance matrix so it fits in GPU constant memory.'' Walk through the arithmetic for $n = 128$ and explain why this argument is not as clean as it sounds.

  \item The \texttt{DIST} macro has parentheses around every argument. Write an example where omitting the parentheses around a single argument would produce a silently wrong result.

  \item At what point in the code does the struct transition from ``partially initialised'' to ``fully valid''? Identify the exact line.

  \item Why does \texttt{tsp\_build\_distance\_matrix} cast to \texttt{(size\_t)} before multiplying \texttt{n * n}? At what value of $n$ would the multiplication overflow a 32-bit signed integer?
\end{enumerate}

\newpage

%% ─────────────────────────────────────────────────────────────────────────────
\section{Implementation Order Summary}
%% ─────────────────────────────────────────────────────────────────────────────

Follow this order. Do not skip steps.

\begin{enumerate}[label=\textbf{Step \arabic*.}, leftmargin=*, itemsep=4pt]
  \item Write \texttt{include/instance.h}: struct, macro, declarations.
  \item Write \texttt{src/instance.c}: \texttt{tsp\_instance\_init} (static), \texttt{tsp\_instance\_free}.
  \item Compile with an empty \texttt{main}. Fix all warnings before proceeding.
  \item Write \texttt{tsp\_instance\_load}: open, scan headers, parse \texttt{DIMENSION}, allocate, find \texttt{NODE\_COORD\_SECTION}, parse coordinates.
  \item Create the test fixtures.
  \item Write \texttt{tests/test\_instance.c}: \texttt{test\_load\_square} only.
  \item Compile and run. Fix failures.
  \item Write \texttt{tsp\_build\_distance\_matrix}.
  \item Add matrix tests (D-01 through D-04). Fix failures.
  \item Add failure-mode tests (F-01, F-02). Fix failures.
  \item Run under Valgrind. Fix all leaks before proceeding.
  \item Write your design notes answering Q1--Q8.
\end{enumerate}

\newpage

%% ─────────────────────────────────────────────────────────────────────────────
\section{Reflection: After Everything Passes}
%% ─────────────────────────────────────────────────────────────────────────────

Once Valgrind is clean and all tests pass, answer these reflection questions. These bridge Phase 1 to the GPU work ahead.

\begin{enumerate}[label=\arabic*., itemsep=6pt]
  \item If the instance size grew to 5000 cities, the full matrix would be $5000^2 \times 8 = 200$\,MB. What options would you have? Name at least two distinct strategies.

  \item GPU global memory is large but has high latency. GPU shared memory is fast but limited ($\approx$48\,KB per thread block on modern hardware). When the GA kernel accesses the distance matrix, what access pattern determines whether shared memory can help?

  \item Your current parser is strict: it only handles one file format subset. If you later need to support TSPLIB files with explicit edge weight matrices (\texttt{EDGE\_WEIGHT\_SECTION}), what part of \texttt{tsp\_instance\_load} would change? What part would stay the same?

  \item \texttt{tsp\_build\_distance\_matrix} and \texttt{tsp\_instance\_load} are currently separate functions. Describe one scenario where that separation would make your Phase 2 implementation significantly easier.

  \item Write one sentence describing why \texttt{free(NULL)} being a no-op is a deliberate design decision in the C standard, not an accident.
\end{enumerate}

%% ─────────────────────────────────────────────────────────────────────────────
\section{Optional Extensions}
%% ─────────────────────────────────────────────────────────────────────────────

Do these only after all tests pass and the design notes are written.

\subsection*{Extension A: Packed symmetric storage}
Store only entries with $i < j$. The index formula for $(i < j)$ in a packed array is:
\[
\text{idx}(i,j) = i \cdot n - \frac{i(i+1)}{2} + j - i - 1
\]
\begin{itemize}
  \item Implement a \texttt{DIST\_PACKED} macro that handles both $i < j$ and $i > j$.
  \item Rerun all tests. Do any break?
  \item Measure: how much less memory is used for $n = 512$?
\end{itemize}

\subsection*{Extension B: \texttt{float} distance matrix}
Switch \texttt{dist} from \texttt{double} to \texttt{float}. Use \texttt{sqrtf} instead of \texttt{sqrt}.
\begin{itemize}
  \item Adjust test tolerances where needed.
  \item For $n = 128$, compare matrix footprint against constant-memory budget.
  \item Answer: does this change improve anything measurable on CPU?
\end{itemize}

\subsection*{Extension C: Command-line loader}
Write a small \texttt{src/main.c} that accepts a \texttt{.tsp} file as a command-line argument, loads it, builds the matrix, and prints the matrix dimensions and the distance from city 0 to city 1.

\begin{lstlisting}[style=shellstyle]
./tsp_info tests/fixtures/square_4.tsp
# Loaded 4 cities.
# dist(0,1) = 1.000000
\end{lstlisting}

\end{document}